\documentclass[10pt]{article}
\usepackage[margin=20mm]{geometry}
\usepackage{amsmath,amssymb,fontspec,unicode-math,xurl,hyperref}
\IfFontExistsTF{DejaVu Serif}{\setmainfont{DejaVu Serif}}{\setmainfont{Latin Modern Roman}}
\setmathfont{latinmodern-math.otf}
\setlength{\parindent}{0pt}\setlength{\parskip}{6pt}\sloppy
\begin{document}
Abstract: A private randomized-response construction contradicts the lower-bound indistinguishability lemma. The upper bounds also apply local regularity assumptions beyond their domain. Valid intermediate calculations and possible repairs are separated from the unestablished conclusions.\par
Author: AI-assisted technical review\par
Date:\par
Title: Technical review: complexity trade-offs in unlearning\par
\# Scope and assessment\par
This review concerns I25-43736, When to Forget? Complexity Trade-offs in Machine Unlearning, in the 17-page PMLR version. Assessment: Major. A private one-step algorithm contradicts the principal indistinguishability lemma used for the lower bound. Separately, the local regularity hypothesis does not justify the unprojected Gaussian-noise upper bounds. These are defects in central derivations; they do not establish that every proposed regime is false under a repaired formulation. The retained-data-only oracle model, exact pretrained optimum and nonatomic-distribution assumption substantially delimit the conclusions.\par
\# A counterexample to Lemma B.2\par
Equation reference: eq:b2. Problem 1 (Major: the stated total-variation bound is false).  Appendix B, Lemma B.2, equation (29), bounds the output distributions for \(g^\gamma\) and \(g^0\) by \[
 \frac{\pi\gamma\sqrt T}{4}
 +\frac{((1-r_f)\gamma-r_f)_+}{2r_f}(e^\epsilon-1+\delta).\] This fails for an algorithm satisfying even pure \(\epsilon\)-unlearning and belonging to the stated first-order algorithm class at \(T=1\).\par
Take \(d=1\), \(L=2\), \(\mu=1\), so \(R=1\). Let \(D_r,D_f,D_f'\) be uniform on \([0,1]\), \([2,3]\), \([4,5]\), respectively. They meet Assumption 2. Use the lemma’s functions \(g^\gamma\), with \(A_\gamma=[0,(1+\gamma)/2]\) on the retain support, and the quadratic loss \(\ell^\gamma(\theta,\xi)=\theta^2/2-g^\gamma(\xi)\theta\). Its derivative has magnitude at most 2 on \([-1,1]\), and it is globally 1-strongly convex. For \(\gamma>r_f/(1-r_f)\) the full optimum is \(\theta^*_\gamma=(1-r_f)\gamma-r_f>0\), while \(\theta^*_0=0\).\par
At the first update, an algorithm may ignore its sampled gradient and output a bit \(Y\in\{0,1\}\). Set \[\Pr(Y=1\mid\theta_0>0)=p_+=\frac{e^\epsilon}{1+e^\epsilon},\qquad
 \Pr(Y=1\mid\theta_0\leq0)=p_-=\frac1{1+e^\epsilon}.\] This is a measurable update of Algorithm 1, using its random seed. For any two initial values and either output atom, the likelihood ratio is at most \(e^\epsilon\). Consequently it satisfies Definition 2 for every loss and every change in the forget distribution, with \(\delta=0\), hence also for the positive \(\delta\) below. It uses neither forget data nor extra oracle information. For longer horizons it can retain the sampled bit in memory; only \(T=1\) is needed here.\par
Choose \(r_f=0.01\), \(\gamma=r_f/(1-r_f)+10^{-6}\), \(\epsilon=1\), \(\delta=10^{-8}\) and \(T=1\). The two output laws are Bernoulli\((p_+)\) and Bernoulli\((p_-)\), with \[d_{\rm TV}=p_+-p_-=\tanh(1/2)\simeq0.462117.\] The claimed bound [eq:b2] (\#eq:b2) is approximately \(0.00802\). The independent scalar calculation records the exact parameters, optimum signs and both probabilities. No simulation or author program is required.\par
The algorithm is not intended to attain small excess risk. That is irrelevant: Lemma B.2 claims an indistinguishability bound for the admissible unlearning algorithms before any utility restriction is applied. The proof’s final equation (32) retains a ceiling \(K'=\lceil((1-r_f)\gamma-r_f)_+/(2r_f)\rceil\), whereas (29) replaces it by the unrounded quantity. An arbitrarily small positive change of the initial optimum can change a private randomized-response distribution by a fixed amount; privacy alone does not provide the fractional-step continuity assumed by that replacement. Other displayed equalities in (31) should also be inequalities with an absolute telescoping difference.\par
Dependency and repair. Theorem 2 explicitly combines Lemmas B.1 and B.2 in (33)–(35); Corollary 1 and the claimed inefficient-region boundary inherit this gap. Restoring a ceiling changes the subsequent condition and does not, by itself, rederive the printed threshold. A new argument, an appropriately restricted algorithm class, or a revised lower-bound range is needed. This counterexample refutes the lemma and the supplied proof, not the theorem independently.\par
\# Local assumptions do not control the noised iterates\par
Problem 2 (Major: the upper-bound argument uses regularity outside its domain).  Assumption 1 imposes \(L\)-Lipschitzness and \(\mu\)-strong convexity only on \(B(0,L/(2\mu))\). Optimization is nevertheless over \(\mathbb R^d\), and Appendix A adds Gaussian noise without projection and applies a bounded-gradient SGD estimate. Neither the noised point nor subsequent iterates are guaranteed to remain in that ball.\par
Even imposing a unique global minimizer does not fix this. With \(L=2\), \(\mu=1\), \(R=1\), consider a data-independent loss \[f_C(\theta)=\tfrac12\theta^2+C(|\theta|-1)_+^4,\qquad C>0.\] It is differentiable and globally 1-strongly convex, is 2-Lipschitz on \([-1,1]\), and has unique minimizer zero for every distribution. Thus it satisfies the paper’s hypothesis, with even stronger convexity. The Appendix A noising rule still has a positive variance \(\sigma^2\) for \(r_f,\kappa>0\). Its excess risk is exactly \[\mathbb E f_C(G)=\sigma^2/2+C\mathbb E(|G|-1)_+^4,\quad G\sim N(0,\sigma^2).\] The last expectation is strictly positive, so the risk is unbounded as \(C\) increases although all prescribed constants remain fixed. For \(r_f=0.1\), \(\epsilon=1\), \(\delta=0.1\) and \(C=100\), independent Gaussian quadrature gives risk \(0.204016\), exceeding Theorem 1’s lower endpoint \(e=0.131036<e_0=0.5\). The failure is of the proposed universal risk estimate, even when the full and retain minimizers are both inside the ball. A data-aware method could observe that this example needs no noise; that does not rescue the specific universal construction used in the proof.\par
The same domain issue affects (22)’s \(\gamma L^2/2\) term and the scratch-SGD argument (39). Restricting the optimization domain and projecting the initial noise and every update would be a plausible repair: projection is postprocessing for privacy, and keeps the regularity hypotheses applicable. Its details, sensitivity calibration and complexity constants must actually be proved. Global strong convexity plus global Lipschitzness on all of \(\mathbb R^d\) cannot simply be substituted.\par
Additional exact-formula discrepancies. With \(q=r_f/(1-r_f)\), Lipschitzness, if valid at the required points, would give a bound proportional to \((\sqrt d\,\kappa q+q^2)L^2/\mu\). The powers of \(q\) in printed (17) are interchanged; (18) also removes denominators and introduces \(e_0=L^2/(8\mu)\) without a valid equality or inequality. At the parameters above (17) is about \(0.555\), whereas (18) is about \(0.117\), so their displayed ordering fails. Equation (25)’s equality drops a factor 64 upon expressing \(L^2R_1^2/e^2\) in terms of \((e_0/e)^2\). The latter factor does not change a big-\(O\) rate, but Theorem 1 states an explicit threshold. Repair of the local-domain issue should include a fresh constant audit.\par
\# Definitions, algorithms and empirical interpretation\par
There is also a formal zero-query inconsistency. Under Algorithm 1 and (5), \(T=0\) returns \(\theta^*\) unchanged, for every update rule; for neighboring forget distributions with distinct full optima this deterministic map cannot be \((\epsilon,\delta)\)-private when \(\delta<1\). Yet Theorem 1 uses free randomized initialization and returns a noised model with zero gradient queries. The cost interpretation reasonably allows this, but the algorithm class must explicitly allow zero-cost initial/final processing. This is a repairable definition mismatch, not an argument against counting only gradient queries.\par
Algorithm 2 also initializes \(\theta_0\) but first accesses \(\theta_1\), uses a decreasing step size and weighted iterates, whereas the proof of Theorem 3 chooses a constant step size optimized for \(T\) in (22)–(24). The proof may support an existence result with a corrected algorithm; it does not establish the displayed algorithm’s bound as written. The experiment’s oracle sensitivity \(\|\theta^*-\theta_r^*\|\) (page 9) uses the retain optimum, and the experiment averages realized threshold-crossing times. Equations (6)–(7) instead take the first deterministic time at which expected risk is small. These quantities need not coincide. Figure 2 is illustrative evidence for chosen procedures, not a measurement of the minimax ratio over all procedures and losses.\par
The non-atomic condition in Assumption 2 is essential for the continuum of Bernoulli biases in the proof. It excludes a fixed finite empirical distribution, despite the motivating finite-dataset example. First-order retain-only access, exact initial training and absence of retained pretraining state beyond the optimum are further substantive scope restrictions. The discussion acknowledges some of these and the unmatched dimension dependence. In particular the asserted exact-unlearning consequence \(\kappa=+\infty\) is outside Theorem 2’s stated \(\delta\geq10^{-8}\) range; a limiting or separate proof is needed.\par
\# Checks and remaining coverage\par
Main pages 1–9 and Appendices A–D (12–17) were read, and Figures 1 and 2 were actually viewed. Lemma B.1’s coupling argument gives the stated two-point squared-distance constant. Lemma C.5’s binomial bound follows from a score-variance bound and integration of \((1-\gamma^2)^{-1/2}\) (with ordering/absolute-value conventions). The expected squared distance in (21) correctly splits into the independent centered-noise variance and squared bias. These valid components do not imply the invalid Lemma B.2 or supply missing domain control. No relevant author repository was acquired; no raw experiment, external citation theorem, historical version or separate inaccessible supplement was independently reproduced. Figure 2’s qualitative regions cannot adjudicate the proof defects.\par
The source SHA-256 is `ec459fbbd69c48ee4ffcecb690deb640a918d797f6872f0aae93169d3a4c2898`. Saved physical-page extracts, independent numerical checks, viewed-page hashes and exact-TeX build records accompany this initial Astra Max review. Scientific adjudication remains pending; human reviews are empty and this is not independent acceptance.\par
\end{document}
